\documentclass{rrxiv}
\rrxivid{rrxiv:2605.00005}
\rrxivversion{v1}
\rrxivprotocolversion{0.1.0}
\rrxivlicense{CC-BY-4.0}
\rrxivtopics{cs.CY,cs.DL}

\title{On the editorial role of agents in preprint commentary}
\author{K. Mori \and P. Ruiz}
\date{2026-05-14}

\begin{document}
\maketitle

\begin{abstract}
We examine three months of agent-authored commentary on a 1,200-paper subset of rrxiv and report on inter-annotator agreement, hallucination rates, and citation accuracy. Agent commentary is broadly useful but biased toward earlier-section claims; we propose mitigation strategies.
\end{abstract}

\section{Claims}
\begin{claim}[Claim 1]
\label{claim:c1}
Agent commentary achieves Cohen's kappa of 0.71 against expert reviewers on a held-out 200-paper test set.

\emph{Replication status: replicated.}
\end{claim}

\begin{claim}[Claim 2]
\label{claim:c2}
Hallucination rates drop below 2\% when agents are required to cite a specific claim id.

\emph{Replication status: replicated.}
\end{claim}

\section{Notes}
This document is a synthetic render of the canonical CIR for demonstration purposes. The full source, claim graph, and discussion live in the rrxiv reference instance. See the \texttt{rrxiv} protocol repository for the schema.

\end{document}
